\section{Graph benchmarks}\label{sec:graphs}

The same question for graphs. We score skeletons only, against the 17-arc Sachs consensus (853 cells, log scale) and against the documented graphs of Causal Chambers recordings: two i.i.d.\ recordings (ten random subsamples of $4000$ shuffled rows each), two further recordings with an effective sample size above $1000$, and seventeen new wind-tunnel and light-tunnel recordings whose actuators are random walks and whose effective sample size, by the heuristic $n(1-\rho)/(1+\rho)$, is about three out of $4000$ rows. For the walk recordings the first $4000$ rows in time order are used. Table~\ref{tab:graphs} gives precision/recall/$F_1$ for $22$ method configurations: the classical rules (PC at two levels, GES, DirectLiNGAM, NOTEARS, DAGMA, PCMCI+), nine published packages at library defaults (experiment 45; AVICI, GOLEM, NOTEARS-MLP, GraN-DAG, DAG-GNN, GRaSP, BOSS, RESIT, CAM-UV), BF-BIC with BOSS, and the methods we developed (PC with a nonparametric test, BTAD, LRG, sequential certificates, Appendix~\ref{app:negative}).

\begin{table}[t]
\caption{Skeleton precision/recall/$F_1$ against the reference (strict), mean over data sets. Sachs: one data set; Chambers iid: $2\times10$ subsamples; ESS$>1000$: $2$ recordings; ESS$<20$: $17$ recordings. A number in parentheses is the number of data sets on which the method finished within $900$\,s when fewer than the full number. Experiments 36, 39, 43, 45.}
\label{tab:graphs}
\scriptsize
\setlength{\tabcolsep}{2pt}
\begin{tabular}{@{}lcccc@{}}
\toprule
Method & Sachs & Chambers iid & ESS $>1000$ & ESS $<20$\\
\midrule
PC $\alpha=$0.05 & 1.00/0.41/0.58 & 0.40/0.38/0.37 & 0.48/0.28/0.35 & 0.36/0.39/0.37\\
PC $\alpha=$0.01 & 1.00/0.35/0.52 & 0.39/0.33/0.34 & 0.51/0.27/0.36 & 0.36/0.36/0.36\\
GES-BIC & 1.00/0.41/0.58 & 0.45/0.64/0.49 & 0.74/0.67/0.70 & 0.26/0.71/0.38\\
DirectLiNGAM & 1.00/0.35/0.52 & 0.55/0.60/0.55 & 0.71/0.63/0.66 & 0.26/0.63/0.37\\
NOTEARS & 1.00/0.18/0.30 & 0.41/0.21/0.27 & 0.33/0.24/0.26 & 0.26/0.45/0.33\\
DAGMA & 1.00/0.29/0.45 & 0.44/0.33/0.36 & 0.44/0.41/0.39 & 0.25/0.52/0.34\\
PCMCI+ & 1.00/0.35/0.52 & 0.39/0.33/0.34 & 0.53/0.29/0.37 & 0.27/0.17/0.20\\
GRaSP & 1.00/0.35/0.52 & 0.53/0.57/0.52 & 0.78/0.63/0.68 & 0.27/0.50/0.35\\
BOSS & 1.00/0.35/0.52 & 0.57/0.61/0.55 & 0.77/0.57/0.65 & 0.29/0.54/0.37\\
BF-BIC + BOSS & 0.90/0.53/0.67 & 0.43/0.63/0.49 & 0.57/0.61/0.58 & 0.30/0.68/0.42\\
NOTEARS-MLP & 0.88/0.41/0.56 & 0.42/0.51/0.43 & 0.45/0.27/0.34 (1) & 0.17/0.25/0.21 (1)\\
GOLEM & 1.00/0.18/0.30 & 0.26/0.39/0.31 & 0.36/0.49/0.41 & 0.25/0.47/0.33\\
DAG-GNN & 1.00/0.18/0.30 & 0.36/0.45/0.36 & 0.45/0.34/0.38 & 0.27/0.45/0.34 (16)\\
GraN-DAG & 1.00/0.06/0.11 & 0.55/0.04/0.06 & 0.20/0.04/0.06 & 0.54/0.07/0.12\\
RESIT & 0.39/0.53/0.45 & 0.33/0.80/0.43 & 0.37/0.74/0.49 & 0.23/0.99/0.37\\
CAM-UV & 1.00/0.06/0.11 & 0.67/0.31/0.42 & 0.74/0.43/0.53 & 0.06/0.01/0.02 (15)\\
AVICI & 0.88/0.41/0.56 & 0.48/0.36/0.37 (14) & 0.54/0.17/0.26 & 0.36/0.23/0.28 (8)\\
PC-GCM $\alpha=$0.05 & 0.89/0.47/0.62 & 0.38/0.33/0.34 & 0.54/0.34/0.41 & 0.35/0.54/0.42\\
PC-GCM $\alpha=$0.01 & 1.00/0.47/0.64 & 0.39/0.30/0.32 & 0.47/0.26/0.33 & 0.34/0.46/0.39\\
BTAD & 0.89/0.47/0.62 & 0.56/0.36/0.41 & 0.82/0.46/0.58 & 0.60/0.16/0.25\\
LRG & 0.90/0.53/0.67 & 0.39/0.48/0.43 & 0.52/0.55/0.53 & 0.51/0.28/0.36\\
CERT-CD nonparametric & 1.00/0.29/0.45 & 0.29/0.16/0.20 & 0.60/0.18/0.28 & 0.24/0.87/0.37\\
\bottomrule
\end{tabular}
\end{table}

\paragraph{The ranking depends on the family.} On the i.i.d.\ chamber recordings the best are DirectLiNGAM and BOSS ($F_1=0.55$) and GRaSP ($0.52$); on the two recordings with a large effective sample size GES ($0.70$), GRaSP ($0.68$) and DirectLiNGAM ($0.66$); on Sachs LRG and BF-BIC with BOSS ($0.67$), PC with the nonparametric test ($0.62$--$0.64$) and BTAD ($0.62$), against $0.58$ for PC and GES. No method has a strict $F_1$ above $0.42$ on the seventeen random-walk recordings, where most methods report tens of false edges. The neural methods are weaker: GraN-DAG finds almost no edges ($F_1\le0.12$), GOLEM and DAG-GNN report $22$ to $36$ false edges per recording, NOTEARS with an MLP timed out on $16$ of $17$ walk recordings ($16$ variables) and AVICI on $9$ of $17$ and on $6$ of $20$ i.i.d.\ subsamples; their means in Table~\ref{tab:graphs} are over the data sets on which they finished and are not comparable to the others.

\paragraph{Resolution and structure of the benchmarks.} One Sachs arc changes recall by $0.06$; $16$ of $22$ configurations have a bootstrap interval that contains the best $F_1$ (Section~\ref{sec:noise}). The chamber graphs are bipartite (actuators to sensors), so skeleton recovery is feature selection from actuators to measurements, and the actuators are recognisable by their marginals (Section~\ref{sec:marg}). An interventional pilot on the light tunnel (a pilot, not pre-specified; \texttt{results/lt\_interventions\_check.json}) shows how little a graph benchmark with documented interventions leaves to a new method: a two-sample test of each sensor against the reference recording under each actuator intervention has precision $1.00$ and recall $0.72$ ($28$ of the $39$ arcs among observed variables, which here are also the ancestral pairs; AUC $0.91$ from the $p$-values), and the missed arcs, all from light sources to their sensors, show no shift in the data ($p$ between $0.06$ and $0.86$); a conditional version (residuals of a boosted regression on the estimated non-descendants) raises recall to $0.79$ at a precision of $0.51$.

\paragraph{What we do not conclude.} The reference networks are small, the Sachs consensus is contested, and several baselines (CAM, SCORE) were not run. The differences between graph methods are mostly within sampling error, and the effective sample size of the walk recordings limits what any method can show. We read the graph results as consistent with the pair results: no rule leads across families.
